\chapter{Geometric limits of generalized Ricci flows}
\label{chap:geometric-limits}

In this chapter we apply the main result of the last section,
bounded curvature at bounded distance, to blow-up limits in order to
establish the existence of a smooth limit for sequences of
generalized Ricci flows. In the first section we establish a blow-up
limit that is defined for some interval of time of positive length,
where the length of the interval of time is allowed to depend on the
limit. In the second section we give conditions under which this
blow-up limit can be extended backwards to make an ancient Ricci
flow. In the third section we construct limits at the singular time
of a generalized Ricci flow satisfying appropriate conditions. We
characterize the ends of the components of these limits. We show
that they are $\epsilon$-horns -- the ends are diffeomorphic to
$S^2\times [0,1)$ and the scalar curvature goes to infinity at the
end. In the fourth section we prove  for any $\delta>0$  that there
are $\delta$-necks sufficiently deep in any $\epsilon$-horn,
provided that the curvature at the other end of the horn is not too
large. Throughout this chapter we fix $\epsilon>0$ sufficiently
small such that all the results of the Appendix hold for $2\epsilon$
and $\alpha=10^{-2}$, and Proposition~\ref{narrows} holds for
$2\epsilon$.


\section{A smooth blow-up limit defined for a small time}


We begin with a theorem that produces a blow-up limit flow defined
on some small time interval.

\begin{theorem}[Existence of a short-time blow-up limit]
\label{thm:short-time-blowup-limit}
\label{smlmtflow}
Fix  canonical neighborhood constants $(C,\epsilon)$, and
non-collapsing constants $r>0,\kappa>0$. Let $({\mathcal
M}_n,G_n,x_n)$ be a sequence of based generalized $3$-dimensional
Ricci flows. We set $t_n={\bf t}(x_n)$ and $Q_n=R(x_n)$. We denote
by $M_n$ the $t_n$ time-slice of ${\mathcal M}_n$. We suppose that:
\begin{enumerate}
\item[(1)] Each $({\mathcal M}_n,G_n)$ either has a time interval of definition contained in $[0,\infty)$
and has  curvature pinched toward positive, or has  non-negative curvature.
\item[(2)] Every point $y_n\in ({\mathcal M}_n,G_n)$ with ${\bf t}(y_n)\le
t_n$ and with $R(y_n)\ge 4R(x_n)$ has a strong
$(C,\epsilon)$-canonical neighborhood.
\item[(3)] $\lim_{n\rightarrow\infty}Q_n=\infty$.
\item[(4)] For each $A<\infty$ the following holds for all $n$ sufficiently large. The ball
$B(x_n,t_n,AQ_n^{-1/2})$ has compact closure in $M_n$ and the flow
is $\kappa$-non-collapsed on scales $\le r$ at each point of
$B(x_n,t_n,AQ_n^{-1/2})$.
\item[(5)] There is $\mu>0$ such that for every $A<\infty$ the following holds
 for all $n$ sufficiently large.
For every $y_n\in B(x_n,t_n,AQ_n^{-1/2})$ the maximal flow line
through $y_n$ extends backwards for a time at least $\mu\left(\mathrm{max}(Q_n,R(y_n))\right)^{-1}$.
\end{enumerate}


Then, after passing to a subsequence and shifting the times of each
of the generalized flows so that $t_n=0$ for every $n$, there is a
geometric limit $(M_\infty,g_\infty,x_\infty)$ of the sequence of
based Riemannian manifolds $(M_n,Q_nG_n(0),x_n)$. This limit is a
complete $3$-dimensional Riemannian manifold of bounded,
non-negative curvature. Furthermore, for some $t_0>0$ which depends
on the curvature bound for $(M_\infty,g_\infty)$ and on $\mu$, there
is a geometric limit Ricci flow defined on $(M_\infty,g_\infty(t)),
-t_0\le t\le 0$, with $g_\infty(0)=g_\infty$.
\end{theorem}

Before beginning the proof of this theorem we establish a lemma that
we shall need both in its proof and also for later applications.


\begin{lemma}[Curvature control from canonical neighborhoods]
\label{lem:curvature-control-canonical-nbhd}
\label{controlnbhd}
Let $({\mathcal M},G)$ be a generalized $3$-dimensional Ricci flow.
Suppose that $r_0>0$ and that any $z\in {\mathcal M}$ with $R(z)\ge
r_0^{-2}$ has a strong $(C,\epsilon)$-canonical
neighborhood. Suppose $z\in
{\mathcal M}$ and ${\bf t}(z)=t_0$. Set
$$r=\frac{1}{2C\sqrt{\max(R(z),r_0^{-2})}}$$ and
$$\Delta t=\frac{1}{16C\left(R(z)+r_0^{-2}\right)}.$$
Suppose that $r'\le r$ and that $|t'-t_0|\le \Delta t$ and let $I$
be the interval with endpoints $t_0$ and $t'$. Suppose that there is
an embedding of $j\colon B(z,t_0,r')\times I$ into ${\mathcal M}$
compatible with time and with the vector field. Then $R(y)\le
2\left(R(z)+r_0^{-2}\right)$ for all $y$ in the image of $j$.
\end{lemma}

\begin{proof}
We first  prove that for any $y\in B(z,t_0,r)$ we have
\begin{equation}\label{Rineq}
R(y)\le \frac{16}{9}(R(z)+r_0^{-2}).\end{equation} Let $\gamma\colon
[0,s_0]\to B(z,t_0,r)$ be a path of length $s_0<r$ connecting
$z=\gamma(0)$ to $y=\gamma(s_0)$. We take $\gamma$ parameterized by
arc length. For any $s\in [0,s_0]$ let $R(s)=R(\gamma(s))$.
According to the strong $(C,\epsilon)$-canonical neighborhood
assumption at any point where $R(s)\ge r_0^{-2}$ we have $|R'(s)|\le
CR^{3/2}(s)$. Let $J\subset [0,s_0]$ be the closed subset consisting
of $s\in [0,s_0]$ for which $R(s)\ge r_0^{-2}$. There are three
possibilities. If $s_0\not\in J$ then $R(y)\le r_0^{-2}$ and we have
established Inequality~(\ref{Rineq}). If $J=[0,s_0]$, then we have
$|R'(s)|\le CR^{3/2}(s)$ for all $s$ in $J$. Using this differential
inequality and the fact that the interval has length at most
$\frac{1}{2C\sqrt{R(z)}}$, we see that $R(y)\le 16R(z)/9$, again
establishing Inequality~(\ref{Rineq}). The last possibility is that
$J\not= [0,s_0]$ but $s_0\in J$. We restrict attention to the
maximal interval of $J$ containing $s_0$. This interval has length
at most $\frac{r_0}{2C}$ and at its initial point $R$ takes the
value $r_0^{-2}$. For every $s$ in this interval by our assumptions
we again have the inequality $|R'(s)|\le CR^{3/2}(s)$, it follows
immediately that $R(y)\le 16r_0^{-2}/9$. This establishes
Inequality~(\ref{Rineq}) in all cases.


Now consider the vertical path $j(\{y\}\times I)$. Let
$R(t)=R(j(y,t))$. Again by the strong canonical neighborhood
assumption $|R'(t)|\le CR^2(t)$ at all points where $R(t)\ge
r_0^{-2}$. Consider the closed subset $K$ of $I$ where $R(t)\ge
r_0^{-2}$. There are three cases to consider: $t'\not\in K$, $t'\in
K\not=I$, or $K=I$.  In the first case, $R(y,t')\le r_0^{-2}$ and we
have established the result. In the second case, let $K'$ be the
maximal subinterval of $K$ containing $t'$. On the interval $K'$ we
have $|R'(t)|\le CR^2(t)$ and at one endpoint $R(t)=r_0^{-2}$. Since
this interval has length at most $r_0^2/16C$, it follows easily that
$R(t')\le 16r_0^{-2}/15$, establishing the result. In the last case
where $K=I$, then, by what we established above, the initial
condition is $R(t_0)=R(y)\le 16(R(z)+r_0^{-2})/9$, and the
differential inequality $|R'(t)|\le CR^2(t)$ holds for all $t\in I$.
Since the length of $I$ is at most $\frac{1}{16C(R(y)+r_0^{-2})}$ we
see directly that $R(t')\le 2(R(z)+r_0^{-2})$, completing the proof
in this case as well.
\end{proof}

Now we begin the proof of Theorem~\ref{thm:short-time-blowup-limit}.



\begin{proof}(of Theorem~\ref{thm:short-time-blowup-limit})
\uses{bcbd,2ndmfdconv,blowupposcurv,partialflowlimit}
We shift the times for the flows so that $t_n=0$ for all $n$. Since
$Q_n$ tends to $\infty$ as $n$ tends to $\infty$, according to
Theorem~\ref{bcbd} for any $A<\infty$, there is a bound
$Q(A)<\infty$ on the scalar curvature of $Q_nG_n(0)$ on
$B_{Q_nG_n}(x_n,0,A)$  for all $n$ sufficiently large. According to
the hypothesis of Theorem~\ref{thm:short-time-blowup-limit}, this means that there is
$t_0(A)>0$ and, for each $n$ sufficiently large, an embedding of
$B_{Q_nG_n}(x_n,0,A)\times [-t_0(A),0]$ into ${\mathcal M}_n$
compatible with time and with the vector field. In fact, we can
choose $t_0(A)$ so that more is true.

\begin{corollary}[Uniform curvature bound on a backward parabolic neighborhood]
\label{cor:uniform-backward-curvature-bound}
\label{cnnbhd2}
\uses{lem:curvature-control-canonical-nbhd}
For each $A<\infty$, let $Q(A)$ be a bound  on the  scalar curvature
of the restriction of $Q_nG_n$ to  $B_{Q_nG_n}(x_n,0,A)$ for all $n$
sufficiently large. Then there exist a constant  $t'_0(A)>0$
depending on $t_0(A)$ and $Q(A)$, and a constant $Q'(A)<\infty$
depending only on $Q(A)$, and, for all $n$ sufficiently large, an
embedding
$$B_{Q_nG_n}(x_n,0,A)\times (-t'_0(A),0]\to{\mathcal
M}_n$$ compatible with time and with the vector field with the
property that  the scalar curvature of the restriction of $Q_nG_n$
to the image of this subset is bounded by $Q'(A)$.
\end{corollary}

\begin{proof}
\uses{lem:curvature-control-canonical-nbhd}
This is immediate from  Lemma~\ref{lem:curvature-control-canonical-nbhd} and Assumption (5)
in the hypothesis of the theorem.
\end{proof}

Now since the curvatures of the $Q_nG_n$ are pinched toward positive
or are non-negative, bounding the scalar curvature above gives a
bound on $|\Rm_{Q_nG_n}|$ on the product
$B_{Q_nG_n}(x_n,0,A)\times (-t'_0(A),0]$. Now we invoke Shi's
theorem (Theorem~\ref{shi}):

\begin{corollary}[Bounded covariant derivatives of curvature along the blow-up sequence]
\label{cor:shi-bound-blowup-sequence}
\uses{shi,cor:uniform-backward-curvature-bound}
For each $A<\infty$ and for each integer $\ell\ge 0$, there is a
constant $C_2$ such that for all $n$ sufficiently large we have
$$|\nabla^\ell\Rm_{Q_nG_n}(x)|\le C_2$$ for all $x\in
B_{Q_nG_n}(x_n,0,A) $.
\end{corollary}


Also, by the curvature bound and  the $\kappa$-non-collapsed
hypothesis we have the following:

\begin{claim}[Non-collapsing volume bound for the blow-up sequence]
\label{claim:volume-noncollapsing-blowup}
There is $\eta>0$ such that for all $n$ sufficiently large
$$\Vol(B_{Q_nG_n}(x_n,0,\eta))\ge \kappa\eta^3.$$
\end{claim}

Now we are in a position to apply Corollary~\ref{2ndmfdconv}. This
implies that, after passing to a subsequence, there is a geometric
limit $(M_\infty,g_\infty,x_\infty)$ of the sequence of based
Riemannian manifolds $(M_n,Q_nG_n(0),x_n)$. The geometric limit is a
complete Riemannian manifold. If the $({\mathcal M}_n,G_n)$ satisfy
the curvature pinched toward positive hypothesis, by
Theorem~\ref{blowupposcurv}, the limit Riemannian manifold
$(M_\infty,g_\infty)$ has non-negative curvature. If the $({\mathcal
M}_n,G_n)$ have non-negative curvature, then it is obvious that the
limit has non-negative curvature. By construction $R(x_\infty)=1$.

In fact, by Proposition~\ref{partialflowlimit} for each $A<\infty$,
there is $t(A)>0$ and, after passing to a subsequence, geometrically
limit flow defined on $B(x_\infty,0,A)\times (-t(A),0]$.


\begin{claim}[Points of large curvature in the limit have canonical neighborhoods]
\label{claim:canonical-neighborhood-persists}
\uses{thm:short-time-blowup-limit}
Any point in $(M_\infty,g_\infty)$ of curvature greater than $4$ has
a $(2C,2\epsilon)$-canonical neighborhood.
\end{claim}


\begin{proof}
\uses{thm:short-time-blowup-limit,canonvary}
The fact that $(M_\infty,g_\infty,x_\infty)$ is the geometric limit
of the $(M_n,Q_nG_n(0),x_n)$ means that we have the following. There
is an exhausting sequence  $V_1\subset V_2\subset\cdots \subset
M_\infty$ of open subsets of $M_\infty$, with compact closure, each
containing $x_\infty$, and for each $n$  an embedding $\varphi_n$ of
$V_n$ into the zero time-slice of ${\mathcal M}_n$ such that
$\varphi_n(x_\infty)=x_n$ and such that the Riemannian metrics
$\varphi_n^*G_n$ converge uniformly on compact sets to $g_\infty$.
Let $q\in M_\infty$ be a point with $R_{g_\infty}(q)>4$. Then for
all $n$ sufficiently large, $q\in V_n$, so that $q_n=\varphi_n(q)$
is defined, and $R_{Q_nG_n}(q_n)>4$. Thus, $q_n$ has an
$(C,\epsilon)$-canonical neighborhood, $U_n$, in ${\mathcal M}_n$;
and, since $R(q_n)>4$ for all $n$,  there is a uniform bound to the
distance from any point of $U_n$ to $q_n$. Thus, there exists $m$
such that for all $n$ sufficiently large $\varphi_n(V_m)$ contains
$U_n$. Clearly as $n$ goes to infinity the Riemannian metrics
$\varphi_n^*(G_n)|_{\varphi_n^{-1}(U_m)}$ converge smoothly to
$g_\infty|_{\varphi_n^{-1}(U_n)}$. Thus, by
Proposition~\ref{canonvary} for all $n$ sufficiently large the
restriction of $g_\infty$ to $\varphi_n^{-1}(U_n)$ contains a
$(2C,2\epsilon)$-canonical neighborhood of $q$.
\end{proof}

\begin{claim}[The short-time blow-up limit has bounded curvature]
\label{claim:limit-bounded-curvature}
The limit Riemannian manifold $(M_\infty,g_\infty)$ has bounded
curvature.
\end{claim}


\begin{proof}
\uses{claim:canonical-neighborhood-persists,localprod,narrows}
First, suppose that $(M_\infty,g_\infty)$ does not have strictly
positive curvature. Suppose that $y\in M_\infty$ has the property
that $\Rm(y)$ has a zero eigenvalue. Fix $A<\infty$ greater
than $d_{g_\infty}(x_\infty,y)$. Then applying
Corollary~\ref{localprod} to the limit flow on
$B(x_\infty,0,A)\times (-t(A),0]$,  we see that the Riemannian
manifold $(B(x_\infty,0,A),g_\infty)$  is locally a Riemannian
product of a compact surface of positive curvature with a
one-manifold. Since this is true for every $A<\infty$ sufficiently
large, the same is true for $(M_\infty,g_\infty)$. Hence
$(M_\infty,g_\infty)$ has a one- or two-sheeted covering that is a
global Riemannian product of a compact surface and one-manifold.
Clearly, in this case the curvature of $(M_\infty,g_\infty)$ is
bounded.


If $M_\infty$ is compact, then it is clear that the curvature is
bounded.

It remains to consider the  case where $(M_\infty,g_\infty)$ is
non-compact and of strictly positive curvature. Since any point of
curvature greater than $4$ has a $(2C,2\epsilon)$-canonical
neighborhood, and since $M_\infty$ is non-compact, it follows that
the only possible canonical neighborhoods for $x\in M_\infty$ are  a
$2\epsilon$-neck centered at $x$  or $(2C,2\epsilon)$-cap whose core
contains $x$.  Each of these canonical neighborhoods contains a
$2\epsilon$-neck. Thus, if $(M_\infty,g_\infty)$ has unbounded and
positive Riemann curvature or equivalently, it has  unbounded scalar
curvature, then it has $(2C,2\epsilon)$-canonical neighborhoods of
arbitrarily small scale, and hence $2\epsilon$-necks of arbitrarily
small scale. But this contradicts Proposition~\ref{narrows}. It
follows from this contradiction that the curvature of
$(M_\infty,g_\infty)$ is bounded.
\end{proof}


To complete the proof of Theorem~\ref{thm:short-time-blowup-limit} it remains to
extend the limit for the $0$ time-slices of the $({\mathcal
M}_n,G_n)$ that we have just constructed to a limit flow defined for
some positive amount of time backward. Since the curvature of
$(M_\infty,g_\infty)$ is bounded, this implies that there is a
bound, $Q$, such that for any $A<\infty$
 the curvature of the restriction of
$Q_nG_n$ to $B_{Q_nG_n}(x_n,0,A)$ is bounded by $Q$ for all $n$
sufficiently large. Thus, we can take the constant $Q(A)$ in
Corollary~\ref{cor:uniform-backward-curvature-bound} to be independent of $A$. According to that
corollary this implies that there is a $t'_0>0$ and $Q'<\infty$ such
that for every $A$ there is an embedding $B_{Q_nG_n}(x_N,0,A)\times
(-t'_0,0]\to {\mathcal M}_n$ compatible with time and with the
vector field so that the scalar curvature of the restriction of
$Q_nG_n$ to the image is bounded by $Q'$ for all $n$ sufficiently
large. This uniform bound on the scalar curvature yields a uniform
bound, uniform in the sense of being independent of $n$, on $|\Rm_{Q_nG_n}|$ on the image of the embedding
$B_{Q_nG_n}(x_N,0,A)\times (-t'_0,0]$.



 Then by Hamilton's result,
Proposition~\ref{partialflowlimit}, we see that, after passing to a
further subsequence, there is a limit flow defined on $(-t_0',0]$.
Of course, the zero time-slice of this limit flow is the limit
$(M_\infty,g_\infty)$. This completes the proof of
Theorem~\ref{thm:short-time-blowup-limit}.
\end{proof}


\section{Long-time blow-up limits}


Now we wish to establish conditions under which we can, after
passing to a further subsequence, establish the existence of a
geometric limit flow defined on $-\infty<t\le 0$. Here is the main
result.


\begin{theorem}[Existence of a long-time (ancient) blow-up limit]
\label{thm:long-time-blowup-limit}
\label{kaplimit}
\uses{thm:short-time-blowup-limit}
Suppose that $\{({\mathcal M}_n,G_n,x_n)\}_{n=1}^\infty$ is a
sequence of generalized $3$-dimensional Ricci flows satisfying all
the hypothesis of Theorem~\ref{thm:short-time-blowup-limit}. Suppose in addition that
there is  $T_0$ with $0<T_0\le \infty$ such that the following
holds. For any $T<T_0$, for each $A<\infty$, and all $n$
sufficiently large, there is an embedding
$B(x_n,t_n,AQ_n^{-1/2})\times (t_n-TQ_n^{-1},t_n]$ into ${\mathcal
M}_n$ compatible with time and with the vector field and at every
point of the image the generalized flow is
$\kappa$-non-collapsed on scales $\le
r$. Then, after shifting the times of the generalized flows so that
$t_n=0$ for all $n$ and passing to a subsequence there is a
geometric limit Ricci flow
$$(M_\infty,g_\infty(t),x_\infty),\ -T_0<t\le 0,$$
for the rescaled generalized flows $(Q_n{\mathcal M}_n,Q_nG_n,x_n)$.
This limit flow is complete and of  non-negative curvature.
Furthermore, the curvature is locally bounded in time. If in
addition $T_0=\infty$, then it is a $\kappa$-solution.
\end{theorem}


\begin{remark}[Comparison with the short-time blow-up limit]
\label{rem:comparison-short-and-long-time-limits}
\uses{thm:short-time-blowup-limit,thm:long-time-blowup-limit}
Let us point out the differences between this result and
Theorem~\ref{thm:short-time-blowup-limit}. The hypotheses of this theorem include all
the hypotheses of Theorem~\ref{thm:short-time-blowup-limit}. The main difference
between the conclusions is that in Theorem~\ref{thm:short-time-blowup-limit} the
amount of backward time for which the limit flow is defined depends
on the curvature bound for the final time-slice of the limit (as
well as how far back the flows in the sequence are defined). This
amount of backward time tends to zero as the curvature of the final
time-slice limit tends to infinity. Here, the amount of backward
time for which the limit flow is defined depends only on how far
backwards the flows in the sequence are defined.
\end{remark}

\begin{proof}
\uses{thm:short-time-blowup-limit,blowupposcurv,narrows,Harnack}
In Theorem~\ref{thm:short-time-blowup-limit} we proved that, after passing to a
subsequence, there is a geometric limit Ricci flow, complete of
bounded non-negative curvature,
$$(M_\infty,g_\infty(t),x_\infty),\ -t_0\le t\le 0,$$
defined for some $t_0>0$. Our next step is to
extend the limit flow all the way back to time $-T_0$.

 \begin{proposition}[Extending a bounded-curvature limit flow backward in time]
\label{prop:extend-limit-flow-backward}
\label{maxT}
\uses{thm:long-time-blowup-limit}
With the notation of, and under the hypotheses of
Theorem~\ref{thm:long-time-blowup-limit}, suppose that there is a geometric limit flow
$(M_\infty,g_\infty(t))$ defined for $-T<t\le 0$ which has
non-negative curvature locally bounded in time. Suppose that
$T<T_0$. Then the curvature of the limit flow is bounded and the
geometric limit flow can be extended to a flow with bounded
curvature defined on $(-(T+\delta),0]$ for some $\delta>0$.
 \end{proposition}


\begin{proof}
\uses{lem:curvature-control-canonical-nbhd,thm:short-time-blowup-limit,cor:curvature-bound-near-basepoint,narrows,Harnack}
 The
argument is by contradiction, so we suppose that there is a $T<T_0$
as in the statement of the proposition. Then the geometric limit
flow on $(-T,0]$ is complete of non-negative curvature and with the
curvature locally bounded in time. First suppose that the scalar
curvature is bounded by, say $Q<\infty$. Fix $T'<T$.
 The Riemannian manifold
$(M_\infty,g_\infty(T'))$ is complete of non-negative curvature with
the scalar curvature, and hence the norm of the Riemann curvature,
bounded by $Q$. Thus, for any $A<\infty$ for all $n$ sufficiently
large, the norm of the Riemann curvature of $Q_nG_n(-T')$ on
$B_{Q_nG_n}(x_n,-T,A)$ is bounded above by $2Q$. Also, arguing as in
the proof of Theorem~\ref{thm:short-time-blowup-limit} we see that any point $y\in
M_\infty$ with $R(y,-T')> 4$ has a $(2C,2\epsilon)$-canonical
neighborhood. Hence, applying Lemma~\ref{lem:curvature-control-canonical-nbhd} as in the
argument in the proof of
 Corollary~\ref{cor:uniform-backward-curvature-bound} shows that for all $n$ sufficiently large,
every point in $B_{Q_nG_n}(x_n,-T',A)$ has a uniform size parabolic
neighborhood on which the Riemann curvature is uniformly bounded,
where both the time interval in the parabolic neighborhood and the
curvature bound on this neighborhood depend only on $C$ and the
curvature bound on $Q$ for the limit flow. According to Hamilton's
result (Proposition~\ref{partialflowlimit}) this implies that, by
passing to a further subsequence, we can extend the limit flow
backward
 beyond $-T'$ a uniform amount of time, say $2\delta$. Taking $T'>T-\delta$
 then gives the desired extension under the condition that the
scalar curvature is bounded on $(-T,0]$.


It remains to show that, provided that $T<T_0$, the scalar curvature
of the limit flow $(M_\infty,g_\infty(t)),\ -T<t\le 0$, is bounded.
 To establish this we need a couple of preliminary results.

\begin{lemma}[Curvature bound on a compact subset from a pointwise minimum bound]
\label{lem:curvature-bound-compact-subset}
\label{compactbd}
Suppose that there is a geometric limit flow defined on $(-T,0]$ for
some $0<T\le T_0$ with $T<\infty$. We suppose that this limit is
complete with non-negative curvature, and with curvature locally
bounded in time. Suppose that $X\subset M_\infty$ is a compact,
connected subset. If $\min_{x\in X}(R_{g_\infty} (x,t))$ is
bounded, independent of $t$, for all $t\in (-T,0]$, then there is a
finite upper bound on $R_{g_\infty}(x,t)$ for all $x\in X$ and all
$t\in(-T,0]$.
\end{lemma}


\begin{proof}
\uses{claim:distance-distortion-bound,bcbd}
Let us begin with:

\begin{claim}[Distance distortion bound under the limit flow]
\label{claim:distance-distortion-bound}
\label{distch}
 Let $Q$ be an upper bound on $R(x,0)$ for all $x\in
M_\infty$. Then for any points $x,y\in M_\infty$ and any $t\in
(-T,0]$ we have
$$d_{t}(x,y)\le d_0(x,y)+16\sqrt{\frac{Q}{3}}T.$$
\end{claim}

\begin{proof}
\uses{cor:distance-integral-bound,Harnack}
Fix $-t_0\in (-T,0]$. Then for any $\epsilon>0$ sufficiently small,
by the Harnack inequality
(the second result in Theorem~\ref{Harnack})
 we have
$$\frac{\partial R}{\partial t}(x,t)\ge -\frac{R(x,t)}{t+T-\epsilon}.$$
Taking the limit as $\epsilon\rightarrow 0$ gives
$$\frac{\partial R}{\partial t}(x,t)\ge -\frac{R(x,t)}{t+T},$$
and hence, fixing $x$,
$$\frac{dR(x,t)}{R(x,t)}\ge \frac{-dt}{(t+T)}.$$
Integrating from $-t_0$ to $0$ shows that
$$\log(R(x,0))-\log(R(x,-t_0))\ge \log(T-t_0)-\log(T),$$
and since $R(x,0)\le Q$, this implies
$$R(x,-t_0)\le Q\frac{T}{T-t_0}.$$
Recalling that $n=3$ and that the curvature is non-negative we see that
$$\Ric(x,-t_0)\le (n-1)\frac{QT}{2}\frac{1}{T-t_0}.$$
Hence by Corollary~\ref{cor:distance-integral-bound}, for all $-t_0\in (-T,0]$ we have
that
$$\dist_{-t_0}(x,y)\le \dist_0(x,y)+8\int_{-t_0}^0\sqrt{\frac{QT}{3(T+t)}}
\le \dist_0(x,y)+16\sqrt{\frac{Q}{3}}T.$$
\end{proof}

It follows immediately from this claim that any compact subset
$X\subset M_\infty$ has uniformly bounded diameter under all the
metrics $g_\infty(t);\ -T<t\le 0$.

By the hypothesis of the lemma there is a constant $C'<\infty$ such
that for each $t\in (-T,0]$ there is $y_{t}\in X$ with
$R_{g_\infty}(y_{t},t)\le C'$. Suppose that the conclusion of the
lemma does not hold. Then there is a sequence $t_m\rightarrow -T$ as
$m\rightarrow\infty$ and points $z_m\in X$ such that
$R_{g_\infty}(z_m,t_m)\rightarrow\infty$ as $m\rightarrow\infty$. In
this case, possibly after redefining the constant $C'$, we can also
assume that there is a point $y_m$ such that $2\le R(y_m,t_m)\le
C'$. Since the sequence $({\mathcal M}_n,Q_nG_n,x_n)$ converges
smoothly to $(M_\infty,g_\infty(t),x_\infty)$ for $t\in (-T,0]$, it
follows that for each $m$ there are sequences $\{y_{m,n}\in
{\mathcal M}_n\}_{n=1}^\infty$ and $\{z_{m,n}\in {\mathcal
M}_n\}_{n=1}^\infty$ with ${\bf t}(y_{m,n})={\bf t}(z_{m,n})=t_m$
converging to $(y_m,t_m)$ and $(z_m,t_m)$ respectively. Thus, for
all $m$ there is $n_0=n_0(m)$ such that for all $n\ge n_0$ we have:
\begin{enumerate}
\item[(1)] $1\le R_{Q_nG_n}(y_{m,n})\le 2C'$,
\item[(2)] $R_{Q_nG_n}(z_{m,n})\ge R_{g_\infty}(z_m,t_m)/2$,
\item[(3)] $d_{Q_nG_n}((y_{m,n}),(z_{m,n}))\le 2\,\operatorname{diam}_{g_\infty(t_m)}(X).$
\end{enumerate}
Because of the third condition and the fact that $X$ has uniformly
bounded diameter under all the metrics $g_\infty(t)$ for $t\in
(-T,0]$, the distance $d_{Q_{m,n}G_{m,n}}(z_{m,n},y_{m,n})$ is
bounded independent of $m$ and $n$ as long as $n\ge n_0$. Because of
the fact that $R_{Q_nG_n}(y_{m,n})\ge 1=R(x_n)$, it follows that any
point $z\in {\mathcal M}_n$ with ${\bf t}(z_m)\le t_{m}$ and with
$R(z)\ge 4R(y_{m,n})$ has a strong $(C,\epsilon)$-canonical
neighborhood. This then
contradicts Theorem~\ref{bcbd} and completes the proof of the lemma.
\end{proof}



Clearly, this argument will be enough to handle the case when
$M_\infty$ is compact. The case when $M_\infty$ is non-compact uses
additional results.

\begin{lemma}[Existence of a point far from a distant point in non-negative curvature]
\label{lem:far-point-with-distant-antipode}
Let $(M,g)$ be a complete, connected, non-compact manifold of
non-negative sectional curvature and let $x_0\in M$ be a point. Then
there is $D>0$, such that for any $y\in M$ with $d(x_0,y)=d\ge D$,
there is $x\in M$ with $d(y,x)=d$ and with $d(x_0,x)>3d/2$.
\end{lemma}


\begin{proof}
\uses{cor:angle-monotonicity}
Suppose that the result is false for $(M,g)$ and $x\in M$. Then
there is a sequence $y_n\in M$ such that setting $d_n=d(x,y_n)$ we
have $\lim_{n\rightarrow\infty}d_n=\infty$ and yet
$B(y_n,d_n)\subset B(x,3d_n/2)$ for every $n$. Let $\gamma_n$ be a
minimal geodesic from $x$ to $y_n$. By passing to a subsequence we
arrange that the $\gamma_n$ converge to a minimal geodesic ray
$\gamma$ from $x$ to infinity in $M$. In particular, the angle at
$x$ between $\gamma_n$ and $\gamma$ tends to zero as
$n\rightarrow\infty$. Let $w_n$ be the point on $\gamma$ at distance
$d_n$ from $x$, and let $\alpha_n=d(y_n,w_n)$. Because $(M,g)$ has
non-negative curvature, by Corollary~\ref{cor:angle-monotonicity}, $\mathrm{lim}_{n\rightarrow\infty}\alpha_n/d_n=0$. In particular, for all $n$
sufficiently large, $\alpha_n<d_n$. This implies that there is a
point $z_n$ on the sub-ray of $\gamma$ with endpoint $w_n$ at
distance $d_n$ from $y_n$. By the triangle inequality,
$d(w_n,z_n)\ge d_n-\alpha_n$. Since $\gamma$ is a minimal geodesic
ray,  $d(z,z_n)=d(z,w_n)+d(w_n,z_n)\ge 2d_n-\alpha_n$. Since
$\alpha_n/d_n\rightarrow 0$ as $n\rightarrow \infty$, it follows
that for all $n$ sufficiently large $d(z,z_n)>3d_n/2$. This
contradiction proves the lemma.
\end{proof}


\begin{claim}[Curvature bound outside a large ball around the base point]
\label{claim:curvature-bound-outside-large-ball}
\label{outerbd}
\uses{lem:far-point-with-distant-antipode}
Fix $D<\infty$ greater than or equal to the constant given in the
previous lemma
 for the Riemannian manifold
$(M_\infty,g_\infty(0))$ and the point $x_\infty$. We also choose
$D\ge 32\sqrt{\frac{Q}{3}}T$. Then for any $y\in M_\infty\setminus
B(x_\infty,0,D)$ the scalar curvature $R_{g_\infty}(y,t)$ is
uniformly bounded for all $t\in (-T,0]$.
\end{claim}

\begin{proof}
\uses{claim:distance-distortion-bound,lem:far-point-with-distant-antipode}
Suppose this does not hold for some $y\in M_\infty\setminus
B(x_\infty,0,D)$. Let $d=d_0(x_\infty,y)$. Of course, $d\ge D$.
Thus, by the lemma there is $z\in M_\infty$ with $d_0(y,z)=d$ and
$d_0(x_\infty,z)>3d/2$. Since the scalar curvature $R(y,t)$ is not
uniformly bounded for all $t\in (-T,0]$, there is $t$ for which
$R(y,t)$ is arbitrarily large and hence $(y,t)$ has an
$(2C,2\epsilon)$-canonical neighborhood of arbitrarily small scale.
By Claim~\ref{claim:distance-distortion-bound}  we have $d_{t}(x_\infty,y)\le
d+8\sqrt{\frac{Q}{3}}T$ and $d_{t}(y,z)\le d+8\sqrt{\frac{Q}{3}}T$.
Of course, since $\Ric\ge 0$ the metric is non-increasing in
time and hence $d\le \min(d_{t}(y,z),d_{t}(x_\infty,y))$ and
$3d/2\le d_{t}(x_0,z)$. Since $y$ has a $(2C,2\epsilon)$-canonical
neighborhood in $(M_\infty,g_\infty(t))$, either $y$ is the center
of an $2\epsilon$-neck in $(M_\infty,g_\infty(t))$ or $y$ is
contained in the core of a $(2C,2\epsilon)$-cap in
$(M_\infty,g_\infty(t))$. (The other two possibilities for canonical
neighborhoods require that $M_\infty$ be compact.)

\begin{claim}[The point is a neck center with geodesics exiting opposite ends]
\label{claim:neck-center-opposite-ends}
$y$ cannot lie in the core of a $(2C,2\epsilon)$-cap in
$(M_\infty,g_\infty(t))$, and hence it is the center of a
$2\epsilon$-neck $N$ in
$(M_\infty,g_\infty(t))$. Furthermore, minimal $g(t)$-geodesics from
$y$ to $x_\infty$ and $z$ exit out of opposite ends of $N$ (see the figure in [Morgan--Tian]).
\end{claim}


\begin{proof}
\uses{claim:distance-distortion-bound}
Let ${\mathcal C}$ be a $(2C,2\epsilon)$-canonical
neighborhood of $y$ in
$(M_\infty,g_\infty(t))$. Since $R(y,t)$ can be arbitrarily large,
we can assume that $d \gg 2CR(y)^{-1/2}$, which is a bound on the
diameter of ${\mathcal C}$. This means that minimal $g(t)$-geodesics
$\gamma_{x_\infty}$ and $\gamma_z$ connecting $y$ to $x_\infty$ and
to $z$, respectively, must exit from ${\mathcal C}$. Let $a$ be a
point on $\gamma_{x_\infty}\cap{\mathcal C}$ close to the complement
of ${\mathcal C}$. Let $b$ be a point at the same $g(t)$-distance
from $y$ on $\gamma_z$. In the case that  ${\mathcal C}$ is a cap
or that it is a $2\epsilon$-neck and $\gamma_{x_\infty}$ and
$\gamma_z$ exit from the same end, then
$d_t(b,y)/d_t(a,y)<4\pi\epsilon$. This means that the angle $\theta$
of the Euclidean triangle with these side lengths at the point
corresponding to $y$ satisfies
$$\cos(\theta)\ge 1-\frac{(4\pi\epsilon)^2}{2}.$$

Recall that $Q$ is the maximum value of $R(x,0)$,  and that by
Claim~\ref{claim:distance-distortion-bound} we have
 $$d\le d_t(x_\infty,y)\le d+16\sqrt{\frac{Q}{3}}T,$$
 with the same inequalities holding with $d_t(z,y)$ replacing $d_t(x_\infty,y)$.
 Also, by construction $d\ge 32\sqrt{\frac{Q}{3}}T$.
We set $a_0=d_t(x_\infty,y)$ and $a_1=d_t(z,y)$. Then by the
Toponogov property we have
$$d_t(x,z)^2\le a_0^2+a_1^2-2a_0a_1\left(
1-\frac{(4\pi\epsilon)^2}{2}\right)=(a_0-a_1)^2+(4\pi\epsilon)^2a_0a_1.$$
Since $|a_0-a_1|\le d/2$ and $a_0,a_1\le 3d/2$ and
$\epsilon<1/8\pi$, it follows that $d_t(x,z)<d$. Since distances do
not increase under the flow, it follows that $d_0(x,z)<d$. This
contradicts the fact that $d_0(x,z)=d$.
\end{proof}


It follows that  the point $y$ is the center of a
$(2C,2\epsilon)$-neck $N$ in $(M_\infty,g_\infty(t))$ and minimal
$g(t)$-geodesics from $y$ to $z$ and to $x_\infty$ exit out of
opposite ends of $N$. This implies that $B(y,t,4\pi R(y,t)^{-1/2})$
separates $x_\infty$ and $z$. Since the curvature of the time-slices
is non-negative, the Ricci flow does not increase distances. Hence,
$B(y,0,4\pi R(y,t)^{-1/2})$ separates $z$ from $x_\infty$. (Notice
that since $d>4\pi R(y,t)^{-1/2}$, neither $z$ nor $x_\infty$ lies
in this ball.) Thus, if  $R(y,t)$ is unbounded as $t\rightarrow -T$
then arbitrarily small $g(0)$-balls centered at $y$ separate $z$ and
$x_\infty$. Since $y$ is distinct from $x_\infty$ and $z$, this is
clearly impossible.
\end{proof}




Next we establish that the curvature near the base point $x_\infty$
is bounded for all $t\in (-T,0]$.

\begin{corollary}[Curvature bound near the base point for all backward times]
\label{cor:curvature-bound-near-basepoint}
\label{ancientlimit}
\uses{lem:curvature-bound-compact-subset,claim:curvature-bound-outside-large-ball}
Suppose that there is a geometric limit flow
$(M_\infty,g_\infty(t))$ of a subsequence defined on $(-T,0]$ for
some $T<\infty$. Suppose that the limit flow is complete of
non-negative curvature with the curvature locally bounded in time.
Then for every $A<\infty$ the scalar curvature $R_{g_\infty}(y,t)$
is uniformly bounded for all $(y,t)\in B(x_\infty,0,A)\times
(-T,0]$.
\end{corollary}

\begin{proof}
\uses{lem:curvature-bound-compact-subset,claim:curvature-bound-outside-large-ball,scalarevol}
First we pass to a subsequence so that a geometric limit flow
$$(M_\infty,g_\infty(t),(x_\infty,0))$$ exists on  $(-T,0]$. We let
$Q$ be the upper bound for $R(x,0)$ for all $x\in M_\infty$. We now
divide the argument into two cases: (i) $M_\infty$ is compact, and
(ii) $M_\infty$ is non-compact.

Suppose that $M_\infty$ is compact. By Proposition~\ref{scalarevol}
we know that
$$\min_{x\in M_\infty}(R_{g_\infty}(x,t))$$
is a non-decreasing function of $t$. Since
$R_{g_\infty}(x_\infty,0)=1$, it follows that for each $t\in
(-T,0]$, we have $\min_{x\in M_\infty}R(x,t)\le 1$, and hence
there is a point $x_{t}\in M_\infty$ with $R(x_{t},t)\le 1$. Now we
can apply Lemma~\ref{lem:curvature-bound-compact-subset} to see that the scalar curvature of
$g_\infty$ is bounded on all of $M_\infty\times (-T,0]$.


If $M_\infty$ is non-compact, choose $D$ as in Lemma~\ref{lem:far-point-with-distant-antipode}.
According to that lemma every point in the boundary of
$B(x_\infty,0,D)$ has bounded curvature under $g_\infty(t)$ for all
$t\in (-T,0]$. In particular, for each $t\in (-T,0]$ the minimum of
$R(x,t)$ over $\overline B(x_\infty,0,D)$ is bounded independent of
$t$. Now apply Lemma~\ref{lem:curvature-bound-compact-subset} to the closure of
$B(x_\infty,0,D)$. We conclude that the curvature of
$B(x_\infty,0,D)$ is uniformly bounded for all $g_\infty(t)$ for all
$t\in (-T,0]$. In particular, $R(x_\infty,t)$ is uniformly bounded
for all $t\in (-T,0]$.


Now for any $A<\infty$ we apply Lemma~\ref{lem:curvature-bound-compact-subset} to the compact
subset $\overline B(x_\infty,0,A)$ to conclude that the curvature is
uniformly bounded on $B(x_\infty,0,A)\times (-T,0]$. This completes
the proof of the corollary.
\end{proof}

Now let us return to the proof of Proposition~\ref{prop:extend-limit-flow-backward}.

\begin{claim}[Uniform curvature bound on an extended backward time interval]
\label{claim:uniform-curvature-bound-extended-interval}
\uses{cor:curvature-bound-near-basepoint}
For each $A<\infty$ and for all $n$ sufficiently large, there are
$\delta>0$ with $\delta\le T_0-T$ and a bound, independent of $n$,
on the scalar curvature of the restriction of $Q_nG_n$ to
$B_{Q_nG_n}(x_n,0,A)\times [-(T+\delta),0]$.
\end{claim}

\begin{proof}
\uses{lem:curvature-control-canonical-nbhd,cor:curvature-bound-near-basepoint}
Fix $A<\infty$ and let $K$ be the bound for the scalar curvature of
$g_\infty$ on $B(x_\infty,0,2A)\times (-T,0]$ from
Corollary~\ref{cor:curvature-bound-near-basepoint}.  Lemma~\ref{lem:curvature-control-canonical-nbhd} shows that
there are $\delta>0$ and a bound in terms of $K$ and $C$ on the
scalar curvature of the restriction of $Q_nG_n$ to
$B_{Q_nG_n}(x_n,0,A)\times [-(T+\delta),0]$.
\end{proof}

Since the scalar curvature is bounded, by the assumption that either
the curvature is pinched toward positive or the Riemann curvature is
non-negative, this implies that the sectional curvatures of $Q_nG_n$
are also uniformly bounded on the products
$B_{Q_nG_n}(x_n,0,A)\times [-(T+\delta),0]$ for all $n$ sufficiently
large. Consequently, it follows that by passing to a further
subsequence we can arrange that the $-T$ time-slices of the
$({\mathcal M}_n,G_n,x_n)$ converge to a limit
$(M_\infty,g_\infty(-T))$.
 This limit manifold
satisfies the hypothesis of Proposition~\ref{narrows} and hence, by
that proposition, it has bounded sectional curvature. This means
that there is a uniform $\delta>0$ such that for all $n$
sufficiently large and for any $A<\infty$ the scalar curvatures (and
hence the Riemann curvatures) of the restriction of $Q_nG_n$ to
$B_{Q_nG_n}(x_n,0,A)\times [-(T+\delta),0]$ are uniformly bounded.
This allows us to pass to a further subsequence for which there is a
geometric limit defined on $(-(T+\delta/2),-T]$. This geometric
limit is complete of  bounded, non-negative curvature. Hence, we
have now constructed a limit flow on $(-(T+\delta/2),0]$ with the
property that for each $t\in (-(T+\delta/2),0]$ the Riemannian
manifold $(M,g(t))$ is complete and of bounded non-negative
curvature. (We still don't know whether the entire flow is of
bounded curvature.) But now invoking Hamilton's Harnack inequality
(Theorem~\ref{Harnack}), we see that the curvature is bounded on
$[-T,0]$. Since we already know it is bounded in $(-T+\delta/2,-T]$,
this completes the proof of the proposition.
\end{proof}

 It follows immediately from Proposition~\ref{prop:extend-limit-flow-backward} that there is a
geometric  limit flow defined on  $(-T_0,0]$. The geometric limit
flow on $(-T_0,0]$ is complete of non-negative curvature, locally
bounded in time.



It remains to prove the last statement in the theorem. So let us
suppose that $T_0=\infty$. We have just established the existence of
a geometric limit flow defined for $t\in (-\infty,0]$.  Since the
$({\mathcal M}_n,G_n)$ either have curvature pinched toward positive
or are of non-negative curvature, it follows from
Theorem~\ref{blowupposcurv} that all time-slices of the limit flow
are complete manifolds of non-negative curvature. Since points of
scalar curvature greater than $4$ have $(2C,2\epsilon)$-canonical
neighborhoods, it follows from Proposition~\ref{narrows} that the
curvature is bounded on each time-slice, and hence universally
bounded by the Harnack inequality (Theorem~\ref{Harnack}). Since for
any $A<\infty$ and every $T<\infty$ the parabolic neighborhoods
$B_{Q_nG_n}(x_n,0,A)\times [-T,0]$ are $\kappa$-non-collapsed on
scales $Q_nr$ for every $n$ sufficiently large,  the limit  is
$\kappa$-non-collapsed on scales $\le \lim_{n\rightarrow
\infty}Q_nr$. Since  $r>0$ and $\lim_{n\rightarrow
\infty}Q_n=\infty$, it follows that the limit flow is
$\kappa$-non-collapsed on all scales. Since $R_{Q_n G_n}(x_n)=1$,
$R_{g_\infty}(x_\infty,0)=1$ and the limit flow is non-flat. This
establishes all the properties need to show that the limit is a
$\kappa$-solution. This completes the proof of
Theorem~\ref{thm:long-time-blowup-limit}.
\end{proof}


\section{Incomplete smooth limits at singular times}

Now we wish to consider smooth limits where we do not blow up, i.e.,
do not rescale the metric. In this case the limits that occur can be
incomplete, but we have strong control over their ends.




\subsection{Assumptions}

We shall assume the following about the generalized Ricci flow
$({\mathcal M},G)$:

\begin{definition}[Standing assumptions for limits at a singular time]
\label{def:standing-assumptions-singular-time}
\label{assumptions}


\begin{enumerate}
\item[(a)] The singular times form a discrete subset of $\Ar$, and
each time slice of the flow at a non-singular time is a compact
$3$-manifold. \item[(b)] The time interval of definition of the
generalized Ricci flow $({\mathcal M},G)$ is contained in
$[0,\infty)$ and its curvature is pinched toward positive.
\item[(c)] There are $r_0>0$ and $C<\infty$, such that any point
$x\in {\mathcal M}$ with $R(x)\ge r_0^{-2}$ has a strong
$(C,\epsilon)$-canonical neighborhood. In particular, for every
$x\in {\mathcal M}$  with $R(x)\ge r_0^{-2}$ the following two
inequalities hold:
$$\left|\frac{\partial R(x)}{\partial t}\right|< CR^2(x),$$
$$|\nabla R(x)|< CR^{3/2}(x).$$
\end{enumerate}
\end{definition}
With these assumptions we can say quite a bit about the limit metric
at time $T$.

\begin{theorem}[Properties of the limit metric at a singular time]
\label{thm:singular-time-limit-properties}
\label{omega}
\uses{def:standing-assumptions-singular-time}
Suppose that $({\mathcal M},G)$ is a generalized Ricci flow defined
for $0\le t<T<\infty$ satisfying the three assumptions given
in~\ref{def:standing-assumptions-singular-time}. Let $T^-<T$ be such that there is a
diffeomorphism $\rho\colon M_{T^-}\times [T^-,T)\to {\bf
t}^{-1}([T^-,T))$ compatible with time and with the vector field.
Set $M=M_{T^-}$ and let $g(t),\ T^-\le t<T$, be the family of
metrics $\rho^*G(t)$ on $M$. Let $\Omega\subset
M$ be the subset of defined by
$$\Omega=\left\{x\in M\bigl|\bigr. \liminf_{t\rightarrow T}R_g(x,t)<\infty\right\}.$$
Then $\Omega\subset M$ is an open subset and there is a Riemannian
metric $g(T)$ with the following properties:
\begin{enumerate}
\item[(1)] As $t\rightarrow T$ the metrics $g(t)|_\Omega$ limit to $g(T)$
 uniformly in the $C^\infty$-topology
on every compact subset of $\Omega$.
\item[(2)] The scalar curvature $R(g(T))$ is a proper function from
$\Omega\to \Ar$ and is bounded below.
\item[(3)] Let
$$\widehat{\mathcal M} = {\mathcal M}\cup_{\Omega\times [T^-,T)}\left(\Omega\times [T^-,T]\right).$$
Then the generalized Ricci flow $({\mathcal M},G)$ extends to a
generalized Ricci flow  $(\widehat {\mathcal M},\widehat G)$.
\item[(4)] Every end of a connected component of
$\Omega$ is contained in a strong
$2\epsilon$-tube.
\item[(5)] Any point $x\in \Omega\times\{T\}$ with $R(x)>r_0^{-2}$
has a strong $(2C,2\epsilon)$-canonical neighborhood in $\widehat{\mathcal M}$.
\end{enumerate}
\end{theorem}

\begin{remark}[Proper functions]
\label{rem:proper-function}
Recall that by definition a function $f$ is proper if the pre-image
under $f$ of every compact set is compact.
\end{remark}





In order to prove this result we establish a sequence of lemmas. The
first in the series establishes that $\Omega$ is an open subset and
also establishes the first two of the above five conclusions.




\begin{lemma}[Omega is open and the metric converges to a limit at time T]
\label{lem:omega-open-limit-metric}
\label{Rproper}
\uses{def:standing-assumptions-singular-time,thm:singular-time-limit-properties}
Suppose that $({\mathcal M},G)$ is a generalized Ricci flow defined
for $0\le t<T<\infty$ satisfying the three assumptions given
in~\ref{def:standing-assumptions-singular-time}. Let $T'<T$ be as in the previous theorem, set
$M=M_{T^-}$, and let $g(t)$ be the family of metrics on $M$ and let
$\Omega\subset M$, each being  as defined in the previous theorem.
Then $\Omega\subset M$ is an open subset of $M$.
 Furthermore, the restriction of the family $g(t)$ to $\Omega$
 converges in the  $C^\infty$-topology, uniformly on compact sets of
 $\Omega$,
 to a Riemannian metric $g(T)$.
Lastly,  $R(g(T))$ is a proper function, bounded below, from
$\Omega$ to $\Ar$.
\end{lemma}

\begin{proof}
\uses{lem:curvature-control-canonical-nbhd,shi}
 We pull back $G$ to $M\times [T^-,T)$ to define a Ricci flow
$(M,g(t)),\ T^{-}\le t<T$. Suppose that $x\in \Omega$. Then there is
a sequence $t_n\rightarrow T$ as $n\rightarrow \infty$ such that
$R(x,t_n)$ is bounded above, independent of $n$, by say $Q$. For all
$n$ sufficiently large we have $T-t_n\le
\frac{1}{16C(Q^2+r_0^{-2})}$. Fix such an $n$. Then, according to
the Lemma~\ref{lem:curvature-control-canonical-nbhd}, there is $r>0$ such that $R(y,t)$ is
uniformly bounded for $y\in B(x,t_n,r)\times [t_n,T)$. This means
that $B(x,t_n,r)\subset \Omega$, proving that $\Omega$ is open in
$M$.

Furthermore, since $R(y,t)$ is bounded on $B(x,t_n,r)\times
[t_n,T)$, it follows from the curvature pinching toward positive
hypothesis that $|\Rm(y,t)|$ is bounded on
$B(x,t_n,r)\times[t_n,T)$. Now applying Theorem~\ref{shi} we see
that in fact $\Rm$ is bounded in the $C^\infty$-topology on
$B(x,t_n,r)\times[(t_n+T)/2 ,T)$. The same is of course also true
for $\Ric$ and hence for $\frac{\partial g}{\partial t}$ in the
$C^\infty$-topology. It then follows that there is a continuous
extension of $g$ to $B(x,t_n,r)\times [t_n,T]$.  Since this is true
for every $x\in \Omega$ we see that $g(t)$ converges in the
$C^\infty$-topology, uniformly on compact subsets of $\Omega$, to
$g(T)$.

Lastly, let us consider the function $R(g(T))$ on $\Omega$. Since
the metric $g(T)$ is a smooth metric on $\Omega(T)$, this is a
smooth function. Clearly, by the curvature pinching toward positive
hypothesis, this function is bounded below. We must show that it is
proper. Since $M$ is compact, it suffices to show that if $x_n$ is a
sequence in $\Omega\subset M$ converging to a point $x\in M\setminus
\Omega$ then $R(x_n,T)$ is unbounded. Suppose that $R(x_n,T)$ is
bounded independent of $n$. It follows from Lemma~\ref{lem:curvature-control-canonical-nbhd}
that there is a positive constant $\Delta t$ such that $R(x_n,t)$ is
uniformly bounded for all $n$ and all $t\in [T-\Delta t,T)$, and
hence, by the same result, there is $r>0$ such that $R(y_n, t)$ is
bounded for all $n$,  all $y_n\in B(x_n,T-\Delta t,r)$, and all
$t\in [T-\Delta t,T)$. Since the $x_n\rightarrow x\in M$, it follows
that for all $n$ sufficiently large that $x\in B(x_n,T-\Delta t,r)$,
and hence $R(x,t)$ is uniformly bounded as $t\rightarrow T$. This
contradicts the fact that $x\not\in \Omega$.
\end{proof}







\begin{definition}[Maximal extension of a generalized Ricci flow to time T]
\label{def:maximal-extension-to-time-T}
\label{hatM}
\uses{thm:singular-time-limit-properties,lem:omega-open-limit-metric}
Let
$$\widehat {\mathcal M}= {\mathcal
M}\cup_{\Omega\times[T^-,T)} \left(\Omega\times
[T^-,T]\right).$$ Since both
${\mathcal M}$ and $\Omega\times [T^-,T]$ have the structure of
space-times and the time functions and vector fields agree on the
overlap, $\widehat {\mathcal M}$ inherits the structure of a
space-time. Let $G'(t),\ T^-\le t\le T$, be the smooth family of
metrics on $\Omega$. The horizontal metrics, $G$, on ${\mathcal M}$
and this family of metrics on $\Omega$ agree on the overlap and
hence define a horizontal metric $\widehat G$ on $\widehat {\mathcal
M}$. Clearly, this metric satisfies the Ricci flow equation, so that
$(\widehat {\mathcal M},\widehat G)$ is a generalized Ricci flow
extending $({\mathcal M},G)$. We call this the {\em maximal
extension of } $({\mathcal M},G)$ to time $T$. Notice that even
though the time-slices $M_t$ of ${\mathcal M}$ are compact, it will
not necessarily be the case that the time-slice $\Omega$ is
complete.
\end{definition}

At this point we have established the first three of the five
conclusions stated in Theorem~\ref{thm:singular-time-limit-properties}. Let us turn to the last
two.

\subsection{Canonical neighborhoods for $(\widehat{\mathcal M},\widehat G)$}

We continue with the notation and assumptions of the previous
subsection. Here we establish the fifth conclusion in
Theorem~\ref{thm:singular-time-limit-properties}, namely the existence of strong canonical
neighborhoods for
$(\widehat{\mathcal M},\widehat G)$


\begin{lemma}[Canonical neighborhoods for points of Omega at time T]
\label{lem:canonical-neighborhoods-at-time-T}
\label{omegacanon}
\uses{def:maximal-extension-to-time-T,def:standing-assumptions-singular-time}
For any $x\in \Omega\times\{T\}$ with $R(x,T)>r_0^{-2}$ one of the
following holds:
\begin{enumerate}
\item[(1)] $(x,T)$ is the center of a strong $2\epsilon$-neck in
$(\widehat{\mathcal M},\widehat G)$. \item[(2)] There is a
$(2C,2\epsilon)$-cap in $(\Omega(T),\widehat G(T))$ whose core
contains $(x,T)$. \item[(3)] There is a $2C$-component of
$\Omega(T)$ that contains $(x,T)$.
\item[(4)] There  is a $2\epsilon$-round component of $\Omega(T)$ that
contains $(x,T)$.
\end{enumerate}
\end{lemma}


\begin{proof}
\uses{lem:curvature-control-canonical-nbhd,lem:omega-open-limit-metric,canonvary}
We fix $x\in \Omega(T)$ with $R(x,T)>r_0^{-2}$. First notice that
 for all $t<T$ sufficiently close to $T$ we
have $R(x,t)>r_0^{-2}$. Thus, for all such $t$ the point $(x,t)$ has
a strong $(C,\epsilon)$-canonical neighborhood in $({\mathcal M},G)\subset (\widehat{\mathcal
M},\widehat G)$. Furthermore, since $\lim_{t\rightarrow
T}R(x,t)=R(x,T)<\infty$, for all $t<T$ sufficiently close to $T$,
there is a constant $D<\infty$ such that for any point $y$ contained
in  a strong $(C,\epsilon)$-canonical neighborhood containing
$(x,t)$, we have $D^{-1}R(x,T)\le R(y,t)\le DR(x,T)$. Again assuming
that $t<T$ is sufficiently close to $T$, by Lemma~\ref{lem:curvature-control-canonical-nbhd}
there is $D'<\infty$ depending only on $D$, $t$, and $r_0$ such that
the curvature $R(y,T)$ satisfies $(D')^{-1}R(x,T)\le R(y,T)\le
D'R(x,T)$. By Lemma~\ref{lem:omega-open-limit-metric} this implies that there is a
compact subset $K\subset \Omega(T)$ containing all the
$(C,\epsilon)$-canonical neighborhoods for $(x,t)$.
 By the same lemma, the metrics $G(t)|_K$
converge uniformly in the $C^\infty$-topology to $G(T)|_K$. If there
is a sequence of $t$ converging to $T$ for which the canonical
neighborhood of $(y,t)$ is an $\epsilon$-round component, resp. a
$C$-component, then $(y,T)$ is contained in a $2\epsilon$-round,
resp. a $2C$-component of $\widehat \Omega$. If there is a sequence
of $t_n$ converging to $T$ so that each $(y,t_n)$ has a canonical
neighborhood ${\mathcal C}_n$ that is a $(C,\epsilon)$-cap whose
core contains $(y,t_n)$, then by Proposition~\ref{canonvary} since
these caps are all contained in a fixed compact subset $K$ and since
the $G(t_n)|_K$ converge uniformly in the $C^\infty$-topology to
$G(T)|_K$, it follows that for any $n$  sufficiently large, the
metric $G(T)$ restricted to ${\mathcal C}_n$ contains a
$(2C,2\epsilon)$-cap ${\mathcal C}$ whose core contains $(y,T)$.



Now we examine the case of strong $\epsilon$-necks.

\begin{claim}[A limit of strong epsilon-necks is a strong 2epsilon-neck]
\label{claim:limit-of-necks-is-neck}
\label{epslimit}
Fix a point $x\in\Omega$. Suppose that there is a sequence
$t_n\rightarrow T$ such that for every $n$, the point $(x,t_n)$ is
the center of a strong $\epsilon$-neck in $\widehat{\mathcal M}$.
Then $(x,T)$ is the center of a strong $2\epsilon$-neck in $\widehat
{\mathcal M}$.
\end{claim}

\begin{proof}
By an overall rescaling we can assume that $R(x,T)=1$. For each $n$
let $N_n\subset \Omega$ and let $\psi_n\colon S^2\times
(-\epsilon^{-1},\epsilon^{-1})\to N_n\times\{t\}$ be a strong
$\epsilon$-neck centered at $(x,t_n)$. Let
$B=B(x,T,2\epsilon^{-1}/3)$. Clearly, for all $n$ sufficiently large
$B\subset N_n$. Thus, for each point $y\in B$ and each $n$ there is
a flow line through $y$ defined on the interval
$(t_n-R(x,t_n)^{-1},t_n]$. Since the $t_n\rightarrow T$ and since
$R(x,t_n)\rightarrow R(x,T)=1$ as $n\rightarrow \infty$, it follows
that there is a flow line through $y$ defined on $(T-1,T]$.

Consider the maps
$$\alpha_n\colon B\times (-1,0]\to \widehat{\mathcal M}$$
that send $(y,t)$ to the value at time $t_n-tR(x,t_n)^{-1}$ of the
flow line through $y$. Pulling back the metric $R(x,t_n)\widehat G$
by $\alpha_n$ produces the restriction of a strong $\epsilon$-neck
structure to $B$. The maps $\alpha_n$ converge uniformly in the
$C^\infty$-topology to the map $\alpha\colon B\times (-1,0]\to
\widehat M$ defined by sending $(y,t)$ to the value of the flowline
through $(y,T)$ at the time $T-t$. Hence, the sequence of metrics
$\alpha^*_n(R(x,t_n))\widehat G$ on $B\times (-1,0]$ converges
uniformly on compact subsets of $B\times (-1,0]$ in the
$C^\infty$-topology to the family $\alpha^*(\widehat G)$.  Then, for
all $n$ sufficiently large, the image $\psi_n(S^2\times
(-\epsilon^{-1}/2,\epsilon^{-1}/2))$ is contained in $B$ and has
compact closure in $B$. Since the family of  metrics
$\psi_n^*\widehat G$ on $B$ converge smoothly to $\psi^*\widehat G$,
it follows that
 for every $n$ sufficiently
large, the restriction of $\psi_n$ to $S^2\times
(-\epsilon^{-1}/2,\epsilon^{-1}/2)$ gives the coordinates showing
that the restriction of the family of metrics $\psi^*(\widehat G)$
to the image $\psi_n(S^2\times (-\epsilon^{-1}/2,\epsilon^{-1}/2))$
is a strong $2\epsilon$-neck at
time $T$.
\end{proof}





This completes the proof of the lemma.
\end{proof}

The lemma tells us that every point  $x\in\Omega\times\{T\}$ with
$R(x)>r_0^{-2}$ has a strong $(2C,2\epsilon)$-canonical
neighborhood. Since, by assumption, points at time before $T$ with
scalar curvature at least $r_0^{-2}$ have strong
$(C,\epsilon)$-canonical neighborhoods, this completes the proof of
the fifth conclusion of Theorem~\ref{thm:singular-time-limit-properties}. It remains to establish
the fourth conclusion of that theorem.

\subsection{The ends of $(\Omega,g(T))$}





\begin{definition}[Strong 2ε-horn]
\label{def:strong-epsilon-horn}
\uses{def:maximal-extension-to-time-T}
 A {\em strong $2\epsilon$-horn} in $(\Omega,g(T))$ is a submanifold of $\Omega$ diffeomorphic
to $S^2\times [0,1)$ with the following properties:
\begin{enumerate}
\item[(1)] The embedding $\psi$ of $S^2\times [0,1)$ into $\Omega$ is a proper map.
\item[(2)] Every point of  the image of this map is the center of a strong
$2\epsilon$-neck in $(\widehat{\mathcal M},\widehat G)$.
\item[(3)] The image of the boundary $S^2\times \{0\}$
is the central sphere of a strong $2\epsilon$-neck.
\end{enumerate}
\end{definition}


\begin{definition}[Strong double 2ε-horn and capped 2ε-horn]
\label{def:strong-double-and-capped-horn}
\uses{def:strong-epsilon-horn}
 A {\em strong double $2\epsilon$-horn} in
$(\Omega,g(T))$ is a component of $\Omega$ diffeomorphic to
$S^2\times (0,1)$ with the property that every point of this
component is the center of a strong $2\epsilon$-neck in $\widehat
{\mathcal M}$. This means that a strong double $2\epsilon$-horn is a
$2\epsilon$-tube and hence is a component of $\Omega$ diffeomorphic
to $S^2\times(-1,1)$. Notice that each end of a strong double
$2\epsilon$-horn contains a strong $2\epsilon$-horn.

 For any $C'<\infty$, a $C'$-{\em capped
$2\epsilon$-horn} in $(\Omega,g(T))$ is a
component of $\Omega$ that is a the union of a the core of a
$(C',2\epsilon)$-cap and a strong $2\epsilon$-horn. Such a component
is diffeomorphic to an open $3$-ball or to a punctured $\Ar P^3$.
\end{definition}

(See the figure in [Morgan--Tian].)


\begin{definition}[The subset $\Omega_\rho$ and horns with boundary contained in it]
\label{def:omega-rho}
\uses{thm:singular-time-limit-properties,def:strong-epsilon-horn}
Fix any $\rho,\ 0<\rho<r_0$. We define $\Omega_\rho\subset \Omega$
to be the closed subset of all $x\in \Omega$ for which $R(x,T)\le
\rho^{-2}$. We say that a strong $2\epsilon$-horn $\psi\colon
S^2\times[0,1)\to \Omega$ has {\em boundary contained in}
$\Omega_\rho$ if its boundary,
$\psi(S^2\times\{0\})$, is contained in $\Omega_\rho$.
\end{definition}

\begin{lemma}[Classification of components of Omega disjoint from $\Omega_\rho$]
\label{lem:omega-component-classification}
\uses{lem:canonical-neighborhoods-at-time-T,def:omega-rho}
Suppose that $0<\rho<r_0$ and that $\Omega^0$ is a component of
$\Omega$ which contains no point of $\Omega_\rho$. Then one of the
following holds:
\begin{enumerate}
\item[(1)] $\Omega^0$ is a strong double $2\epsilon$-horn and is diffeomorphic to $S^2\times \Ar$.
\item[(2)] $\Omega^0$ is a  $2C$-capped $2\epsilon$-horn and is diffeomorphic to $\Ar^3$
or to a punctured $\Ar P^3$.
\item[(3)] $\Omega^0$ is a compact component and is the union of the cores of
two $(2C,2\epsilon)$-caps and a strong $2\epsilon$-tube. It is
diffeomorphic to $S^3$, $\Ar P^3$ or $\Ar P^3\#\Ar P^3$.
\item[(4)] $\Omega^0$ is a compact $2\epsilon$-round component and is diffeomorphic to a compact
manifold of constant positive curvature.
\item[(5)] $\Omega^0$ is a compact component that fibers over $S^1$
with fibers $S^2$.
\item[(6)] $\Omega^0$ is a compact
$2C$-component and is diffeomorphic to $S^3$ or to $\Ar P^3$.
\end{enumerate}
\end{lemma}



\begin{proof}
\uses{lem:canonical-neighborhoods-at-time-T,epstopology}
Let $\Omega^0$ be a component of $\Omega$ containing no point of
$\Omega_\rho$. Then for every  $x\in \Omega^0$, we have
$R(x,T)>r_0^{-2}$.   Therefore, by Lemma~\ref{lem:canonical-neighborhoods-at-time-T} $(x,T)$
has a $(2C,2\epsilon)$-canonical neighborhood. Of course, this
entire neighborhood is contained in $\widehat{\mathcal M}$ and hence
is contained in $\Omega^0$ (or, more precisely, in the case of
strong $2\epsilon$-necks in the union of maximum backward flow lines
ending at points of $\Omega^0$). If the canonical neighborhood of
$(x,T)\in \Omega^0$ is a $2C$-component or is an $2\epsilon$-round
component, then of course $\Omega^0$ is that $2C$-component or
$2\epsilon$-round component. Otherwise, each point of $\Omega^0$ is
either the center of a strong $2\epsilon$-neck or is contained in
the core of a $(2C,2\epsilon)$-cap. We have chosen $2\epsilon$
sufficiently  small so that the result follows from
Proposition~\ref{epstopology}.
\end{proof}



\begin{remark}[Possibly infinitely many double horns]
\label{rem:possibly-infinitely-many-horns}
\uses{def:strong-double-and-capped-horn}
We do not claim that there are only finitely many such components;
in particular, as far as we know there may be infinitely double
$2\epsilon$-horns.
\end{remark}

It follows immediately from this lemma that if $X$ is a component of $\Omega$
not containing any point of $\Omega_\rho$, then every end of $X$ is contained
in a strong $2\epsilon$-tube. To complete the proof of Theorem~\ref{thm:singular-time-limit-properties}, it
remains only to establish the same result for the components of $\Omega$ that
meet $\Omega_\rho$. That is part of the content of the next lemma.


\begin{lemma}[Finitely many horns meet $\Omega_\rho$]
\label{lem:finitely-many-horns-meeting-omega-rho}
\label{2epshorns}
\uses{def:omega-rho,def:standing-assumptions-singular-time}
Let $({\mathcal M},G)$ be a generalized $3$-dimensional Ricci flow
defined for $0\le t<T<\infty$ satisfying
Assumptions~\ref{def:standing-assumptions-singular-time}. Fix $0<\rho<r_0$.  Let
$\Omega^0(\rho)$ be the union of all components of $\Omega$
containing points of $\Omega_\rho$. Then $\Omega^0(\rho)$ has
finitely many components and is a union of a compact set and
finitely many strong $2\epsilon$-horns each of which is disjoint
from $\Omega_\rho$ and has its boundary contained in
$\Omega_{\rho/2C}$.
\end{lemma}

\begin{proof}
\uses{lem:canonical-neighborhoods-at-time-T,lem:omega-component-classification,def:omega-rho,Xcontainedin}
Since $R\colon \Omega\times\{T\}\to \Ar$ is a proper function
bounded below, $\Omega_\rho$ is compact. Hence, there are only
finitely many components of $\Omega$ containing points of
$\Omega_\rho$. Let $\Omega^0$ be a non-compact component of $\Omega$
containing a point of $\Omega_\rho$, and let ${\mathcal E}$ be an
end of $\Omega^0$. Let
$$X=\{x\in \Omega^0\bigl|\bigr.
R(x)\ge 2C^2\rho^{-2}\}.$$ Then $X$ is a closed set and contains a
neighborhood of the end ${\mathcal E}$. Since $\Omega^0$ contains a
point of $\Omega_\rho$, $\Omega^0\setminus X$ is non-empty. Let
$X_0$ be the connected component of $X$ that contains a neighborhood
of ${\mathcal E}$. This is a closed, connected set every point of
which has  a $(2C,2\epsilon)$-canonical neighborhood. Since $X_0$
includes an end of $\Omega^0$, no point of $X_0$ can be contained in
an $\epsilon$-round component nor in a $C$-component. Hence, every
point of $X_0$ is either the center of a strong $2\epsilon$-neck or
is contained in the core of a $(2C,2\epsilon)$-cap. Since
$2\epsilon$ is sufficiently small to invoke
Proposition~\ref{Xcontainedin}, the latter implies that $X_0$ is
contained either in a $2\epsilon$-tube which is a union of strong
$2\epsilon$-necks centered at points of $X_0$ or $X_0$ is contained
in a $2C$-capped $2\epsilon$-tube where the core of the cap contains
a point of $X_0$. ($X_0$ cannot be contained in a double capped
$2\epsilon$-tube since the latter is compact.) In the second case,
since this capped tube contains an end of $\Omega^0$, it is in fact
equal to $\Omega^0$. Since a point of $X_0$ is contained in the core
of the $(2C,2\epsilon)$-cap, the curvature of this point is at most
$2C^2\rho^{-2}$ and hence the curvature at any point of the cap is
at least $2C\rho^{-2}>\rho^{-2}$. This implies that the cap is
disjoint from $\Omega_{\rho}$. Of course, any $2\epsilon$-neck
centered at a point of $X_0$ has curvature at least $C^2\rho^{-2}$
and hence is also disjoint from $\Omega_\rho$. Hence, if $\Omega^0$
is a $2C$-capped $2\epsilon$-tube and there is a point of $X_0$ in
the core of the cap, then this component is disjoint from
$\Omega_\rho$, which is a contradiction. Thus,  $X_0$ is contained
in a $2\epsilon$-tube made up of strong $2\epsilon$-necks centered
at points of $X_0$.



This proves that $X_0$ is contained in a strong $2\epsilon$-tube,
$Y$, every point of which has curvature $\ge C^2\rho^{-2}$. Since
$X_0$ is closed but not the entire component $\Omega^0$, it follows
that $X_0$ has a frontier point $y$. Of course,
$R(y)=2C^2\rho^{-2}$. Let $N$ be the strong $2\epsilon$-neck
centered at $y$ and let $S_N^2$ be its central $2$-sphere. Clearly,
every $y'\in S^2_N$ satisfies $R(y')\le 4C^2\rho^{-2}$, so that
$S^2_N$ is contained in $\Omega_{\rho/2C}$. Let $Y'\subset Y$ be the
complementary component of $S^2_N$ in $Y$ that contains a
neighborhood of the end ${\mathcal E}$. Then the closure of $Y'$ is
the required strong $2\epsilon$-horn containing a neighborhood of
${\mathcal E}$, disjoint from $\Omega_\rho$ and with boundary
contained in $\Omega_{\rho/2C}$.


 The last thing to see is
that there are only finitely many such ends in a given component
$\Omega^0$. First suppose that the boundary $2$-sphere of one of the
$2\epsilon$-horns is homotopically trivial in $\Omega^0$. Then this
$2$-sphere separates $\Omega^0$ into two components one of which is
compact and hence $\Omega^0$ has only one boundary component. Thus,
we can assume that all the boundary $2$-spheres of the
$2\epsilon$-horns are homotopically non-trivial. Suppose that two of
these $2\epsilon$-horns containing different ends of $\Omega^0$ have
non-empty intersection. Let $N$ be the $2\epsilon$-neck whose
central $2$-sphere is the boundary of one of the $2\epsilon$-horns.
Then the boundary of the other $2\epsilon$-horn is also contained in
$N$. This means that
 the
union of the two $2\epsilon$-horns and $N$ is a component of
$\Omega$. Clearly, this component has exactly two
ends. Thus, we can assume that all the $2\epsilon$-horns with
boundary in $\Omega_{\rho/2C}$ are disjoint. If two of the
$2\epsilon$-horns have boundary components that are topologically
parallel in $\Omega^0\cap \Omega_{\rho/2C}$ (meaning that they are
the boundary components of a compact submanifold diffeomorphic to
$S^2\times I$), then $\Omega^0$ is diffeomorphic to $S^2\times
(0,1)$ and has only two ends.
 By compactness of
$\Omega_{\rho/2C}$, there can only be finitely many disjoint
$2\epsilon$-horns with non-parallel, homotopically noon-trivial
boundaries in $\Omega^0\cap \Omega_{\rho/2C}$. This completes the
proof of the fact that each component of $\Omega_{\rho/2C}$ has only
finitely many ends.
\end{proof}


This completes the proof of Theorem~\ref{thm:singular-time-limit-properties}.


\section{Existence of strong $\delta$-necks sufficiently deep in
a $2\epsilon$-horn}


We keep the notation and assumptions of the previous section.

\begin{theorem}[Strong $\delta$-necks sufficiently deep in a 2ε-horn]
\label{thm:deep-delta-necks-in-horn}
\label{hexist}
\uses{def:omega-rho,def:standing-assumptions-singular-time}
Fix $\rho>0$. Then for any $\delta>0$ there is an
$0<h=h(\delta,\rho)\le \min(\rho\cdot\delta,\rho/2C)$,
implicitly depending on $r$ and $(C,\epsilon)$ which are fixed, such
that for any generalized Ricci flow $({\mathcal M},G)$ defined for
$0\le t<T<\infty$ satisfying Assumptions~\ref{def:standing-assumptions-singular-time} and for
any $2\epsilon$-horn ${\mathcal H}$ of $(\Omega,g(T))$ with boundary
contained in $\Omega_{\rho/2C}$, every point $x\in {\mathcal H}$
with $R(x,T)\ge h^{-2}$ is the center of a strong
$\delta$-neck in $(\widehat{\mathcal
M},\widehat G)$ contained in ${\mathcal H}$. Furthermore, there is a
point $y\in {\mathcal H}$ with $R(y)=h^{-2}$ with the property that
the central $2$-sphere of the $\delta$-neck centered at $y$ cuts off
an end of the ${\mathcal H}$ disjoint from $\Omega_\rho$ (see the figure in [Morgan--Tian]).
\end{theorem}




\begin{proof}
\uses{thm:short-time-blowup-limit,thm:splitting-theorem,blowupposcurv,localprod,rm>0,2DGSS}
The proof of the first statement is by contradiction. Fix $\rho>0$
and $\delta>0$ and suppose that there is no $0<h\le \mathrm{min}(\rho\cdot\delta,\rho/2C)$ as required. Then there is a sequence
of generalized Ricci flows $({\mathcal M}_n,G_n)$ defined for $0\le
t<T_n<\infty$ satisfying Assumptions~\ref{def:standing-assumptions-singular-time} and points
$x_n\in {\mathcal M}_n$ with ${\bf t}(x_n)=T_n$ contained in
$2\epsilon$-horns ${\mathcal H}_n$ in $\Omega_n$ with boundary
contained in $(\Omega_n)_{\rho/2C}$ with $Q_n=R(x_n)\rightarrow
\infty$ as $n\rightarrow \infty$ but such that no $x_n$ is the
center of a strong $\delta$-neck in $({\mathcal M}_n, G_n)$. Form
the maximal extensions, $(\widehat{\mathcal M}_n,\widehat G_n)$, to
time $T$ of the $({\mathcal M}_n,G_n)$.

\begin{claim}[The rescaled sequence satisfies the short-time blow-up hypotheses]
\label{claim:rescaled-sequence-satisfies-hypotheses}
\label{hyp}
\uses{thm:short-time-blowup-limit}
The sequence $(\widehat{\mathcal M}_n,\widehat G_n,x_n)$ satisfies
the five hypothesis of Theorem~\ref{thm:short-time-blowup-limit}.
\end{claim}

\begin{proof}
\uses{thm:short-time-blowup-limit}
By our assumptions, Hypotheses (1) and (3) of
Theorem~\ref{thm:short-time-blowup-limit} hold for this sequence.
 Also, we are assuming that any point $y\in {\mathcal M}_n$ with
 $R(y)\ge r_0^{-2}$ has a strong $(C,\epsilon)$-canonical neighborhood. Since
$R(x_n)=Q_n\rightarrow \infty$ as $n\rightarrow\infty$ this means
that for all $n$ sufficiently large, any point $y\in \widehat
{\mathcal M}_n$ with $R(y)\ge R(x_n)$ has a strong
$(C,\epsilon)$-canonical neighborhood. This establishes Hypothesis
(2) in the statement of Theorem~\ref{thm:short-time-blowup-limit}.

Next, we have:

\begin{claim}[Balls around $x_n$ eventually lie in the horn with compact closure]
\label{claim:ball-contained-in-horn-compact-closure}
For any $A<\infty$ for all $n$ sufficiently large,
$B(x_n,0,AQ_n^{-1/2})$ is contained in the $2\epsilon$-horn
${\mathcal H}_n$ and has compact closure in ${\mathcal M}_n$.
\end{claim}

\begin{proof}
\uses{bcbd}
Any point  $z\in\partial {\mathcal H}_n$ has scalar curvature at
most $16C^2\rho^{-2}$ and there is a $2\epsilon$-neck centered at
$z$. This means that for all $y$ with
$d_{G_n}(z,y)<\epsilon^{-1}\rho/2C$ we have $R(y)\le
32C^2\rho^{-2}$. Hence, for all $n$ sufficiently large,
$d_{G_n}(x_n,z)>\epsilon^{-1}\rho/2C$, and thus
$d_{Q_nG_n}(x_n,z)>Q^{1/2}_n\epsilon^{-1}\rho/2C$. This implies
that, given $A<\infty$, for all $n$ sufficiently large, $z\not\in
B_{Q_nG_n}(x_n,0,A)$. Since this is true for all $z\in \partial
{\mathcal H}_n$, it follows that for all $n$ sufficiently large
$B_{Q_n,G_n}(x_n,0,A)\subset {\mathcal H}_n$. Next, we must show
that, for all $n$ sufficiently large, this ball has compact closure.
That is to say, we must show that for every $A$ for all $n$
sufficiently large the distance from $x_n$ to the end of the horn
${\mathcal H}_n$ is greater than $AQ_n^{-1/2}$. If not, then since
the curvature at the end of ${\mathcal H}_n$ goes to infinity for
each $n$, this sequence would violate Theorem~\ref{bcbd}.
\end{proof}

Because $B(x_n,0,AQ_n^{-1/2})$ is contained in a $2\epsilon$-horn,
it is $\kappa$-non-collapsed on scales $\le r$ for a universal
$\kappa>0$ and $r>0$. Also, because every point in the horn is the
center of a strong $2\epsilon$-neck, for every $n$ sufficiently
large and every $y\in B(x_n,0,AQ_n^{-1/2})$ the flow is defined on
an interval through $y$ defined for backward time $R(y)^{-1}$.


This completes the proof that all the hypotheses of
Theorem~\ref{thm:short-time-blowup-limit} hold and establishes Claim~\ref{claim:rescaled-sequence-satisfies-hypotheses}.
\end{proof}



We form a new sequence of generalized Ricci flows from the
$({\mathcal M}_n, G_n)$ by translating by $-T_n$, so that the final
time-slice is at ${\bf t}_n=0$, where ${\bf t}_n$ is the time
function for ${\mathcal M}_n$.







Theorem~\ref{thm:short-time-blowup-limit} implies that, after passing to a
subsequence, there is a limit flow
$(M_\infty,g_\infty(t),(x_\infty,0)),\ t\in[-t_0,0])$ defined for
some $t_0>0$ for the  sequence $(Q_n\widehat{\mathcal
M}_n,Q_n\widehat G_n,x_n)$. Because of the curvature pinching toward
positive assumption, by Theorem~\ref{blowupposcurv}, the limit Ricci
flow has non-negative sectional curvature. Of course,
$R(x_\infty)=1$ so that the limit $(M_\infty,g_\infty(0))$ is
non-flat.

\begin{claim}[The blow-up limit splits as a round $S^2$ times $\mathbb{R}$]
\label{claim:limit-splits-as-cylinder}
$(M_\infty,g_\infty(0))$ is isometric to the product $(S^2,h)\times
(\Ar,ds^2)$, where $h$ is a metric of non-negative curvature on
$S^2$ and $ds^2$ is the usual Euclidean metric on the real line.
\end{claim}

\begin{proof}
\uses{thm:splitting-theorem}
Because of the fact that the $({\mathcal M}_n,G_n)$
have curvature pinched toward positive, and since $Q_n$ tend to
$\infty$ as $n$ tends to infinity, it follows that the geometric
limit $(M_\infty,g_\infty)$ has non-negative curvature. In
${\mathcal H}_n$ take a minimizing geodesic ray $\alpha_n$ from
$x_n$ to the end of ${\mathcal H}_n$ and a minimizing geodesic
$\beta_n$ from $x_n$ to $\partial {\mathcal H}_n$. As we have seen,
the lengths of both $\alpha_n$ and $\beta_n$ tend to $\infty$ as
$n\rightarrow \infty$. By passing to a subsequence, we can assume
that the $\alpha_n$ converge to a minimizing geodesic ray $\alpha$
in $(M_\infty,g_\infty)$ and that the $\beta_n$ converge to a
minimizing geodesic ray $\beta$ in $(M_\infty,g_\infty)$. Since, for
all $n$, the union of $\alpha_n$ and $\beta_n$ forms a piecewise
smooth ray in ${\mathcal H_n}$ meeting the central $2$-sphere of a
$2\epsilon$-neck centered at $x_n$ in a single point and at this
point crossing from one side of this $2$-sphere to the other, the
union of $\alpha$ and $\beta$ forms a proper, piecewise smooth map
of $\Ar$ to $M_\infty$ that meets  the central $2$-sphere of a
$2\epsilon$-neck centered at $x_\infty$ in a single point and
crosses from one side to the other at the point. This means that
$M_\infty$ has at least two ends. Since $(M_\infty,g_\infty)$ has
non-negative curvature, according to Theorem~\ref{thm:splitting-theorem}, this
implies that $M_\infty$ is a product of a surface with $\Ar$. Since
$M$ has non-negative curvature, the surface has non-negative
curvature. Since $M$ has positive curvature at at least one point,
the surface is diffeomorphic to the $2$-sphere.
\end{proof}


According to Theorem~\ref{thm:short-time-blowup-limit}, after passing to a subsequence
there is a limit flow defined on some interval of the form
$[-t_0,0]$ for $t_0>0$. Suppose that, after passing to a subsequence
there is a limit flow defined on $[-T,0]$ for some $0<T<\infty$. It
follows that for any $t\in [-T,0]$, the Riemannian manifold
$(M_\infty,g_\infty(t))$ is of non-negative curvature and has two
ends. Again by Theorem~\ref{thm:splitting-theorem}, this implies that for every
$t\in [-T,0]$ the Riemannian manifold $(M_\infty,g_\infty(t))$ is a
Riemannian product of a metric of non-negative curvature on $S^2$
with $\Ar$. Thus, by Corollary~\ref{localprod} the Ricci flow is a
product of a Ricci flow  $(S^2,h(t))$ with the trivial flow on
$(\Ar,ds^2)$. It now follows from Corollary~\ref{rm>0} that for
every $t\in (-T,0]$ the curvature of $g_\infty(t)$ on $S^2$ is
positive.


Let $M_n$ be the zero time-slice of ${\mathcal M}_n$. Since
$(M_\infty,g_\infty,x_\infty)$ is the geometric limit of the
$(M_n,Q_nG_n(0),x_n)$, there is an exhausting sequence $x_\infty\in
V_1\subset V_2\subset\cdots$ of open subsets of $M_\infty$ with
compact closure and embeddings $\varphi_n\colon V_n\to M_n$ sending
$x_\infty$ to $x_n$ such that $\varphi_n^*(Q_nG_n(0))$ converges in
the $C^\infty$-topology, uniformly on compact sets, to $g_\infty$.

\begin{claim}[Backward flow lines with controlled curvature near the cylinder limit]
\label{claim:backward-flow-line-curvature-bound}
For any $z\in M_\infty$ for all $n$ sufficiently large, $z\in V_n$,
so that $\varphi_n(z)$ is defined. Furthermore, for all $n$
sufficiently large, there is a backward flow line through
$\varphi_n(z)$ in the generalized Ricci flow $(Q_n{\mathcal
M}_n,Q_nG_n)$ defined on the interval
$(-T-(R_{Q_nG_n}^{-1}(\varphi_n(z),0)/2),0]$. The scalar curvature
is bounded above on this entire flow line by $R(\varphi_n(z),0)$.
\end{claim}


\begin{proof}
\uses{neckcurv,defncanonnbhd}
Of course, for any compact subset $K\subset M_\infty$ and any $t'<T$
for all $n$ sufficiently large, $K\subset V_n$, and there is an
embedding $\varphi_n(K)\times [-t',0]\subset Q_n{\mathcal M}_n$
compatible with time and the vector field. The map $\varphi_n$
defines a map $Q_n^{-1}\varphi_n\colon K\times [-Q_n^{-1}t',0]\to
{\mathcal M}_n$. Since the scalar curvature of the limit is
positive, and hence bounded away from zero on the compact set
$K\times [-t',0]$ and since $Q_n\rightarrow \infty$ as $n$ tends to
infinity the following is true: For any compact subset $K\subset M$
and any $t'<T$, for all $n$ sufficiently large, the scalar curvature
of $G_n$ on the image $Q_n^{-1}\varphi_n(K)\times [-Q_n^{-1}t',0]$
is greater than $r_0^{-2}$, and hence for all $n$ sufficiently
large, every point in $Q_n^{-1}\varphi_n(K)\times [-Q_n^{-1}t',0]$
has a strong $(C,\epsilon)$-canonical neighborhood in ${\mathcal
M}_n$. Since having a strong $(C,\epsilon)$-canonical neighborhood
is invariant under rescaling, it follows that for all $n$
sufficiently large, every point of $\varphi_n(K)\times [-t',0]$ has
a strong $(C,\epsilon)$-canonical neighborhood.

Next we claim that, for all $n$ sufficiently large and for any $t\in
[-t',0]$, the point $(\varphi_n(z),t)$ is the center of a strong
$\epsilon$-neck. We have already seen that for all $n$ sufficiently
large $(\varphi_n(z),t)$ has a strong $(C,\epsilon)$-canonical
neighborhood. Of course, since $M_\infty$ is non-compact, for $n$
sufficiently large, the canonical neighborhood of $(\varphi_n(z),t)$
must either be a $(C,\epsilon)$-cap or a strong $\epsilon$-neck. We
shall rule out the possibility of a $(C,\epsilon)$-cap, at least for
all $n$ sufficiently large.

To do this,  take $K$ to be a neighborhood of  $(z,0)$ in the limit
$(M_\infty,g_\infty(0))$ with the topology of $S^2\times I$ and with
the metric being the product of a positively curved metric on $S^2$
with the Euclidean metric on $I$. We take $K$ to be sufficiently
large to contain the $2C$-ball centered at $(z,0)$. Because the
limit flow is the product of a positively curved flow on $S^2$ with
the trivial flow on $\Ar$, the flow is distance decreasing. Thus,
for every $t\in [-t',0]$ the submanifold $K\times\{t\}$ contains the
ball in $(M_\infty,g_\infty(t))$ centered at $(z,t)$ of radius $2C$.
 For every $n$
sufficiently large, consider the submanifolds $\varphi_n(K)\times
\{t\}$ of $(M_n,Q_nG_n(t))$.  Since the metrics
$\varphi_n^*Q_nG_n(t)$ are converging uniformly for all $t\in
[-t',0]$ to the product flow on $K$, for all $n$ sufficiently large
and any $t\in [-t',0]$, this submanifold contains the $C$-ball
centered at $(\varphi_n(z),t)$ in $(M_n,Q_nG_n(t))$. Furthermore,
the maximal curvature two-plane at any point of
$\varphi_n(K)\times\{t\}$ is almost tangent to the $S^2$-direction
of $K$. Hence, by Lemma~\ref{neckcurv} the central $2$-sphere of any
$\epsilon$-neck contained $\varphi_n(K)\times\{t\}$ is almost
parallel to the $S^2$-factors in the product structure on $K$ at
every point. This implies that the central $2$-sphere of any such
$\epsilon$-neck is isotopic to the $S^2$-factor of
$\varphi_n(K)\times\{t\}$. Suppose that $(\varphi_n(z),t)$ is
contained in the core of a $(C,\epsilon)$-cap ${\mathcal C}$. Then
${\mathcal C}$ is contained in $\varphi_n(K)\times\{t\}$. Consider
the $\epsilon$-neck $N\subset {\mathcal C}$ that is the complement
of the core of ${\mathcal C}$. Its central $2$-sphere, $\Sigma$, is
isotopic in $K$ to the $2$-sphere factor of $K$, but this is absurd
since $\Sigma$ bounds a $3$-ball in the ${\mathcal C}$. This
contradiction shows that for all $n$ sufficiently large and all
$t\in [-t',0]$, it is not possible for $(\varphi_n(z),t)$ to be
contained in the core of a $(C,\epsilon)$-cap. The only other
possibility is that for all $n$ sufficiently large and all $t\in
[-t',0]$ the point $(\varphi_n(z),t)$ is the center of a strong
$\epsilon$-neck in $(M_n,Q_nG_n(t))$.

Fix $n$ sufficiently large. Since, for all $t\in [-t',0]$, the point
$(\varphi_n(z),t)$ is the center of a strong $\epsilon$-neck, it
follows from Definition~\ref{defncanonnbhd} that for all $t\in
[-t',0]$ we have $R(\varphi_n(z),t)\le R(\varphi_n(z),0)$ (this
follows from the fact that the partial derivative in the
time-direction of the scalar curvature of a strong $\epsilon$-neck
of scale one is positive and bounded away from $0$). It also follows
from Definition~\ref{defncanonnbhd} that the flow near
$(\varphi_n(z),-t')$ extends backwards to time
$$-t'-R^{-1}_{Q_nG_n}(\varphi_n(z,t'))<-t'-R^{-1}_{Q_nG_n}(\varphi_n(z,0)),$$
with the same inequality for scalar curvature holding for all $t$ in
this extended interval. Applying this for $t'<T$ but sufficiently
close to $T$ establishes the last statement in the claim, and
completes the proof of the claim.
\end{proof}

Let $Q_0$ be the upper bound of the scalar curvature of
$(M_\infty,g_\infty(0))$.  By the previous claim, $Q_0$ is also an
upper bound for the curvature of $(M_\infty,g_\infty(-t'))$ for any
$t'<T$. Applying Theorem~\ref{thm:short-time-blowup-limit} to the flows
$(M_n,Q_nG_n(t)),\, -t'-Q_0^{-1}/2< t\le -t'$, we conclude that
there is $t_0$ depending only on the bound of the scalar curvature
of $(M_\infty,g_\infty(-t'))$, and hence depending only on $Q_0$,
such that, after passing to a further subsequence the limit flow
exists for $t\in [-t'-t_0,-t']$. Since the limit flow already exists
on $[-t',0]$, we conclude that, for this further subsequence, the
limit flow exists on $[-t'-t_0,0]$. Now apply this with
$t'=T-t_0/2$. This proves that if, after passing to a subsequence,
there is a limit flow defined on $[-T,0]$, then, after passing to a
further subsequence there is a limit flow defined on $[-T-t_0/2,0]$
where $t_0$ depends only on $Q_0$, and in particular, is independent
of $T$. Repeating this argument with $T+(t_0/2)$ replacing $T$, we
pass to a further subsequence so that the limit flow is defined on
$[-T-t_0,0]$. Repeating this inductively, we can find a sequence of
subsequences so that for the $n$ subsequence the limit flow is
defined on $[-T-nt_0,0]$. Taking a diagonal subsequence produces a
subsequence for which the limit is defined on $(-\infty,0]$.


The limit flow is the product of a flow on $S^2$ of positive
curvature defined for $t\in (-\infty,0]$ and the trivial flow on
$\Ar$. Now, invoking Hamilton's result (Corollary~\ref{2DGSS}), we
see that the ancient solution of positive curvature on $S^2$ must be
a shrinking round $S^2$. This means that the limit flow is the
product of the shrinking round $S^2$ with $\Ar$, and implies that
for all $n$ sufficiently large there is a strong $\delta$-neck
centered at $x_n$. This contradiction proves the existence of $h$ as
required.


Now let us establish the last statement in Theorem~\ref{thm:deep-delta-necks-in-horn}. The
subset of ${\mathcal H}$ consisting of all $z\in {\mathcal H}$ with
$R(z)\le \rho^{-2}$ is compact (since $R$ is a proper function), and
disjoint from any $\delta$-neck of scale $h$ since $h<\rho/2C$. On
the other hand, for any point $z\in {\mathcal H}$ with $R(z)\le
\rho^{-2}$ take a minimal geodesic from $z$ to the end of ${\mathcal
H}$. There must be a point $y$ on this geodesic with $R(y)=h^{-2}$.
The $\delta$-neck centered at $y$ is disjoint from $z$ (since
$h<\rho/2C$) and hence this neck separates $z$ from the end of
${\mathcal H}$. It now follows easily that there is a point $y\in
{\mathcal H}$ with $R(y)=h^{-2}$ and such that the central
$2$-sphere of the $\delta$-neck centered at $y$ divides ${\mathcal
H}$ into two pieces with the non-compact piece disjoint from
$\Omega_\rho$.
\end{proof}


\begin{corollary}[Monotone choice of $h(\rho,\delta)$]
\label{cor:monotone-h-function}
\label{hmonotone}
\uses{thm:deep-delta-necks-in-horn}
We can take the function $h(\rho,\delta)$ in the last lemma to be
$\le \delta\rho$, to be a weakly monotone non-decreasing function of
$\delta$ when $\rho$ is fixed, and to be a weakly monotone
non-decreasing function of $\rho$ when $\delta$ is held fixed.
\end{corollary}


\begin{proof}
\uses{thm:deep-delta-necks-in-horn}
If $h$ satisfies the conclusion of Theorem~\ref{thm:deep-delta-necks-in-horn} for $\rho$
and $\delta$ and if $\rho'\ge \rho$ and $\delta'\ge \delta$ then $h$
also satisfies the conclusion of Theorem~\ref{thm:deep-delta-necks-in-horn} for $\rho'$
and $\delta'$. Also, any $h'\le h$ also satisfies the conclusion of
Theorem~\ref{thm:deep-delta-necks-in-horn} for $\delta$ and $\rho$. Take a sequence
$(\delta_n,\rho_n)$ where each of the sequences $\{\delta_n\}$ and
$\{\rho_n\}$ is a monotone decreasing sequence with limit $0$. Then
we choose $h_n=h(\rho_n,\delta_n)\le \rho_n\delta_n$ as in the
statement of Theorem~\ref{thm:deep-delta-necks-in-horn}. We of course can assume that
$\{h_n\}_n$ is a non-increasing sequence of positive numbers with
limit $0$. Then for any $(\rho,\delta)$ we take the largest $n$ such
that $\rho\ge \rho_n$ and $\delta\ge \delta_n$, and we define
$h(\rho,\delta)$ to be $h_n$ for this value of $n$. This constructs
the function $h(\delta,\rho)$ as claimed in the corollary.
\end{proof}
